% =========================================================
% Proof: Interior Extremum Theorem
% Source: volume-iii/analysis/differentiation/notes/notes-extrema.tex
% Mode: standard
% =========================================================

\subsection*{Interior Extremum Theorem}
\label{prf:interior-extremum-theorem}

\begin{remark*}[Return]
\hyperref[thm:interior-extremum-theorem]{$\leftarrow$ Back to Theorem (Interior Extremum Theorem) in Notes}
\end{remark*}

% ---------------------------------------------------------
% Proof I — Learning Version
% ---------------------------------------------------------
\subsubsection*{Proof I \textemdash\ Learning Version (Contradiction)}
\label{prf:interior-extremum-theorem-naive}

\begin{proof}
Let $I \subseteq \mathbb{R}$ be an open interval, let $f : I \to \mathbb{R}$,
and let $c \in I$. Assume $f$ has a relative maximum at $c$ and is
differentiable at $c$. The claim is that $f'(c) = 0$.

Suppose for contradiction that $f'(c) \neq 0$.

By \ByDef{Relative Maximum} there exists $\delta > 0$ such that
$(c - \delta, c + \delta) \subseteq I$ and $f(x) - f(c) \leq 0$ for all
$x \in (c - \delta, c + \delta)$.

\ProofStep{1}{Case $f'(c) > 0$.}
Since $f'(c) = \lim_{x \to c} \frac{f(x)-f(c)}{x-c} > 0$, by
\ByThm{Sign Preservation for Limits} there exists a sub-interval
$(c, c + \delta')$ on which the difference quotient is positive. For
$x \in (c, c + \delta')$ the displacement $x - c > 0$, so positivity of
the quotient requires $f(x) - f(c) > 0$, i.e.\ $f(x) > f(c)$. This
contradicts $f(x) - f(c) \leq 0$ on $(c - \delta, c + \delta)$.

\ProofStep{2}{Case $f'(c) < 0$.}
Since $f'(c) < 0$, by \ByThm{Sign Preservation for Limits} there exists
a sub-interval $(c - \delta', c)$ on which the difference quotient is
negative. For $x \in (c - \delta', c)$ the displacement $x - c < 0$, so
negativity of the quotient requires $f(x) - f(c) > 0$, i.e.\ $f(x) > f(c)$.
This again contradicts $f(x) - f(c) \leq 0$.

\ProofStep{3}{Conclusion.}
Both cases lead to a contradiction, so $f'(c) \neq 0$ is false.
Therefore $f'(c) = 0$.

The case of a relative minimum follows by applying the same argument to
$-f$: $-f$ has a relative maximum at $c$, so $(-f)'(c) = 0$, giving
$f'(c) = 0$.
\end{proof}

% ---------------------------------------------------------
% Proof II — Professional Standard
% ---------------------------------------------------------
\subsubsection*{Proof II \textemdash\ Professional Standard (Direct)}
\label{prf:interior-extremum-theorem-direct}

\begin{proof}
Let $I \subseteq \mathbb{R}$ be an open interval, let $f : I \to \mathbb{R}$,
and let $c \in I$. Assume $f$ has a relative maximum at $c$ and is
differentiable at $c$. The claim is that $f'(c) = 0$.

Since $f$ has a relative maximum at $c$, by \ByDef{Relative Maximum} there
exists $\delta > 0$ such that $(c - \delta, c + \delta) \subseteq I$ and
\[
  f(x) - f(c) \leq 0
  \qquad \text{for all } x \in (c - \delta, c + \delta).
\]

\ProofStep{1}{Sign of the difference quotient from the right.}
For $x \in (c, c + \delta)$, the increment $f(x) - f(c) \leq 0$ and the
displacement $x - c > 0$, so
\[
  \frac{f(x) - f(c)}{x - c} \leq 0.
\]
Passing to the right-hand limit and applying \ByThm{Order Preservation for
Limits},
\[
  f'_+(c) = \lim_{x \to c^+} \frac{f(x) - f(c)}{x - c} \leq 0.
\]

\ProofStep{2}{Sign of the difference quotient from the left.}
For $x \in (c - \delta, c)$, the increment $f(x) - f(c) \leq 0$ and the
displacement $x - c < 0$, so
\[
  \frac{f(x) - f(c)}{x - c} \geq 0.
\]
Passing to the left-hand limit and applying \ByThm{Order Preservation for
Limits},
\[
  f'_-(c) = \lim_{x \to c^-} \frac{f(x) - f(c)}{x - c} \geq 0.
\]

\ProofStep{3}{Conclude $f'(c) = 0$.}
Since $f$ is differentiable at $c$, by \ByThm{Differentiability and
One-Sided Derivatives}, $f'(c) = f'_+(c) = f'_-(c)$. From Steps 1 and 2,
\[
  0 \geq f'(c) \geq 0,
\]
therefore $f'(c) = 0$.

The case of a relative minimum follows by applying the above to $-f$ at $c$.
\end{proof}

% ---------------------------------------------------------
% Shared remark blocks
% ---------------------------------------------------------
\begin{remark*}[Proof shape]
The direct proof reads the sign of the difference quotient on each side of
$c$ separately. From the right, a non-positive increment over a positive
displacement gives a non-positive quotient, so $f'_+(c) \leq 0$. From the
left, a non-positive increment over a negative displacement gives a
non-negative quotient, so $f'_-(c) \geq 0$. Differentiability forces the
two one-sided derivatives to agree, pinning $f'(c) = 0$. No contradiction
is needed.

\FlashProof{1. Relative maximum gives $f(x) - f(c) \leq 0$ near $c$.
  2. For $x > c$: quotient $\leq 0$, so $f'_+(c) \leq 0$ by order
  preservation.
  3. For $x < c$: quotient $\geq 0$, so $f'_-(c) \geq 0$ by order
  preservation.
  4. Differentiability gives $f'_+(c) = f'_-(c) = f'(c)$, so $f'(c) = 0$.}
\end{remark*}

\begin{remark*}[Self-assessment of Proof I]
Proof I is correct. At three months of proof writing, producing a valid
contradiction argument with proper case analysis and correct citation of
a named limit theorem is genuinely good work. The following notes identify
specific patterns to grow out of.

\textbf{Inefficiency: contradiction when a direct argument exists.}
Proof I opens with "suppose for contradiction that $f'(c) \neq 0$" and then
cases on the sign of $f'(c)$. But the conclusion $f'(c) = 0$ can be reached
without contradiction at all: the sign inequalities on the one-sided
derivatives directly force $0 \geq f'(c) \geq 0$. Contradiction is a valid
strategy but it is often chosen by default when a direct argument is
available. The habit to build: before assuming the negation, ask whether
the conclusion can be squeezed out of the hypotheses directly. Here it can.

\textbf{Naivety: two cases for what is one phenomenon.}
Cases 1 and 2 in Proof I treat $f'(c) > 0$ and $f'(c) < 0$ separately, each
producing a slightly different contradiction. They look like two different
arguments, but they are the same argument — the difference quotient has a
sign mismatch with the maximum condition — applied once to the right and once
to the left. Proof II recognises this as a single phenomenon (split by the
sign of $x - c$) and handles both sides in two parallel steps rather than
two cases. When two cases in a proof look nearly identical, that is a signal
that there is a unifying observation waiting to be named.

\textbf{Extra theorem required: Sign Preservation vs.\ Order Preservation.}
Proof I invokes Sign Preservation for Limits, which is a theorem about the
behaviour of a function in a neighbourhood given the sign of its limit.
Proof II invokes Order Preservation for Limits, which is a theorem about the
behaviour of a limit given the sign of the function on a neighbourhood.
These are converses of each other. In this proof, the information flows from
the function to the limit (the maximum condition gives the sign of the
quotient, which is then passed to the limit), so Order Preservation is the
natural tool. Sign Preservation runs in the opposite direction and requires
an extra step — finding the sub-neighbourhood — before the contradiction
can be assembled. Recognising the direction of information flow in a proof
is a skill that eliminates unnecessary theorem invocations.

\textbf{What Proof I does well.}
The case structure is clear and explicit. The contradiction is properly
closed in each case. The minimum case is handled by reduction to $-f$,
which is the correct move and was stated cleanly. The hypothesis is fully
set up before the case split. These are all marks of a well-organised proof,
and they carry forward into the professional version unchanged.
\end{remark*}

\begin{remark*}[Commentary]
\textbf{(a) Why does the maximum give $f(x) - f(c) \leq 0$?}
The definition of a relative maximum says $f(c)$ is the largest value in
some neighbourhood of $c$. Every other value in that neighbourhood is at
most $f(c)$, which is exactly $f(x) - f(c) \leq 0$. This single inequality
is the only fact about the maximum used anywhere in either proof.

\textbf{(b) Why split into left and right in Proof II?}
The numerator $f(x) - f(c)$ is always $\leq 0$ near $c$. But the denominator
$x - c$ changes sign: positive for $x > c$, negative for $x < c$. Splitting
the analysis by side lets the sign of the quotient be determined cleanly on
each half, yielding two one-sided inequalities that together pin $f'(c)$.

\textbf{(c) Why does Order Preservation apply?}
Order Preservation for Limits says: if a function is $\leq 0$ on a full
deleted half-neighbourhood, its one-sided limit (when it exists) satisfies
the same inequality. The difference quotient is $\leq 0$ for all $x \in
(c, c+\delta)$, not just eventually, so the hypothesis is met directly.

\textbf{(d) Why does differentiability pin $f'(c) = 0$?}
Differentiability at $c$ forces the one-sided derivatives to exist and
agree: $f'_+(c) = f'_-(c) = f'(c)$. Step 1 gives $f'(c) \leq 0$ and Step 2
gives $f'(c) \geq 0$. The only real number satisfying both simultaneously
is $0$.
\end{remark*}

\begin{remark*}[Dependencies]
\begin{itemize}
  \item \hyperref[def:relative-maximum]{Definition of Relative Maximum}:
    Supplies the neighbourhood and the inequality $f(x) - f(c) \leq 0$
    (both proofs).
  \item \hyperref[thm:order-preservation-limits]{Order Preservation for
    Limits}: Passes sign inequalities through the one-sided limits
    (Proof II, Steps 1 and 2).
  \item \hyperref[thm:sign-preservation-limits]{Sign Preservation for
    Limits}: Guarantees a sub-interval of definite sign from the limit
    (Proof I, Steps 1 and 2).
  \item \hyperref[thm:differentiability-one-sided]{Differentiability and
    One-Sided Derivatives}: Forces $f'_+(c) = f'_-(c) = f'(c)$
    (Proof II, Step 3).
\end{itemize}

\FlashUses{def:relative-maximum, thm:order-preservation-limits,
thm:differentiability-one-sided}
\end{remark*}